\documentclass[11pt]{article}
\usepackage{amsmath,amssymb,amsthm}
\usepackage[margin=1in]{geometry}
\newtheorem{theorem}{Theorem}
\newtheorem{lemma}{Lemma}
\begin{document}
\title{V111: Target-Compatible Minimum-Overlap Flows for Ternary MUX Avoidance}
\author{NC0\_k-Avoid Laboratory}
\date{}
\maketitle

A signed essential ternary MUX gate has the form
\[
 g=o\oplus \operatorname{MUX}(s\oplus p_s,a\oplus p_a,b\oplus p_b).
\]
Fixing its output gives two exact binary clauses.  Selecting one branch yields an
implication whose source phase is fixed by the selector polarity and whose arrival
phase is chosen by the output target.

\begin{lemma}[Arbitrary compatible overlap]
Let two directed branch cycles start at one selector $v$ with opposite first
source phases.  The cycles may share output gates.  If every output gate used by
both cycles receives the same required target bit from both traversals, then the
combined target extends to a word outside the circuit range.
\end{lemma}

\begin{proof}
Along each cycle choose targets so every selected implication arrives at the
source phase of the next selected branch, except the last, which arrives at the
complement of the initial phase.  A cycle beginning with $\alpha$ therefore
forces $x_v=1\oplus\alpha$.  The two initial phases are opposite, so the two
cycles force opposite values of $x_v$.  Shared output gates cause no target
conflict by hypothesis; unused MUX clauses only strengthen the fixed-output
2-CNF.
\end{proof}

For a repeated selector $v$ and opposite-phase first branches with destinations
$d_0,d_1$, remove outputs whose selector is $v$.  Split each remaining output
$h$ with two parallel unit-capacity arcs
\[
 h_{\rm in}\to h_{\rm out}
\]
of costs $0$ and $1$.  All other network arcs have cost zero and capacity two.
A super-source sends one unit to each $d_i$ and the sink is $v$.

\begin{lemma}[Overlap accounting]
For an integral two-flow in this network, the total cost equals the number of
output gates used by both decoded source-to-sink routes.
\end{lemma}

\begin{proof}
The first unit through a gate uses its zero-cost copy.  A second unit can pass
only through the cost-one copy.  Thus a gate contributes one exactly when both
units use it.
\end{proof}

\begin{theorem}[Recognizable minimum-overlap class]
There is a deterministic polynomial-time procedure that scans repeated selectors
and opposite-phase first branches, computes a minimum-cost integral two-flow,
and constructs a missing output whenever the decoded routes are target-compatible
on every shared output gate.
\end{theorem}

\begin{proof}
Only two flow units are sent.  Successive shortest augmenting paths in the
residual network compute a minimum-cost integral flow in polynomial time; reverse
arcs are handled by a negative-edge-capable shortest-path routine.  Decode the
two source-labeled routes, compare the target requirement at each shared gate,
and apply the first lemma when all requirements agree.
\end{proof}

The theorem is a sound recognizer.  It does not assert that the particular
minimum-cost decomposition chosen by the implementation finds a compatible pair
whenever some other equal-cost or higher-cost compatible pair exists.

\begin{theorem}[Nested-bottleneck separation]
For every $d\ge 1$ there is an exact-stretch essential MUX family with
\[
 n=5(d+1),\qquad m=n+1,
\]
whose two canonical return routes have minimum shared-gate overlap exactly $d$
and compatible targets on all $d$ shared gates.  For $d\ge2$, the V110
single-shared-gate procedure still returns a nested bottleneck, whereas V111
constructs a missing output in polynomial time.
\end{theorem}

\begin{proof}
Use $d+1$ two-lobe layers connected by $d$ hub MUX gates.  Every return path from
either central destination to the root selector must traverse every hub layer,
and each hub has exactly one outgoing MUX output.  Hence both routes must use all
$d$ shared hub gates.  Canonical left/right routes share no other output, proving
minimum overlap $d$.  Their phase propagation assigns the same target to each
shared hub gate, so the first theorem applies.  If only the first bottleneck is
upgraded as in V110 and $d\ge2$, the next hub remains unit-capacity and yields a
nested bottleneck.
\end{proof}

\begin{theorem}[Exact V102 backdoor]
On the depth-$d$ family,
\[
 \beta_{\rm V102}=3(d+1)=\frac{3n}{5}.
\]
\end{theorem}

\begin{proof}
For one two-variable lobe exiting through hub $z$, the two local V102 constraints
require at least one lobe variable when $z$ is selected and both lobe variables
when $z$ is not selected.  Each layer has two lobes, so its lobe contribution is
at least $4-2z$.  Summing over the $d+1$ exit hubs gives
\[
 \sum_j(4-2z_j)+\sum_j z_j
 =4(d+1)-\sum_j z_j\ge3(d+1).
\]
Selecting every exit hub and one variable from each lobe meets all constraints,
so equality holds.
\end{proof}

\paragraph{Boundary.}
V111 does not decide target compatibility over all equal-cost flow decompositions
or higher-cost reroutings.  It does not prove all essential MUX/bijunctive
circuits are in P, unrestricted $NC^0_3$-Avoid, a new general circuit lower bound,
or P versus NP.  Novelty and priority are not established.

\end{document}
